
\chapter{TRIỂN KHAI THỰC TẾ, THÁCH THỨC VÀ XU HƯỚNG TƯƠNG LAI}
\label{chap:chuong_5}
% ==========================================
% SECTION 5.1 (NÂNG CẤP VỚI CASE STUDY THỰC TẾ)
% ==========================================
\section{Kịch bản ứng dụng và Nghiên cứu tình huống thực tế (Case Studies)}
\label{sec:5_1}

\subsection{Lực lượng lao động từ xa và làm việc kết hợp: Case Study Google BeyondCorp}
\label{subsec:5_1_1}

Trong môi trường làm việc kết hợp (\textit{Hybrid Work}), nhân viên thường xuyên truy cập tài nguyên nội bộ từ mạng gia đình hoặc mạng công cộng không đáng tin cậy. Điển hình thực tế rõ nét nhất cho kịch bản này là dự án \textbf{BeyondCorp} do Google triển khai từ năm 2011 đến năm 2017 \cite{ward2014}.

\begin{itemize}
    \item \textbf{Bối cảnh và vấn đề gặp phải:} Sau chiến dịch tấn công có chủ đích \textit{Operation Aurora} (cuối năm 2009) nhắm vào hạ tầng mạng doanh nghiệp, Google nhận thấy mô hình bảo mật chu vi dựa trên tường lửa biên và VPN truyền thống tồn tại điểm yếu chí mạng: một khi thiết bị vượt qua cổng VPN, hệ thống mặc định tin cậy toàn bộ lưu lượng xuất phát từ dải IP nội bộ, tạo điều kiện cho kẻ tấn công di chuyển sang ngang. Bên cạnh đó, việc duy trì kết nối VPN cho hơn 100.000 nhân viên toàn cầu gây suy giảm hiệu năng và tăng tải vận hành.
    \item \textbf{Mô hình kiến trúc ZTNA áp dụng:} Google loại bỏ hoàn toàn hệ thống VPN truy cập từ xa và xây dựng kiến trúc Ủy nhiệm nhận biết định danh (\textit{Identity-Aware Proxy}) theo mô hình \textit{Client-initiated}:
    \begin{enumerate}
        \item \textit{Kho quản lý thiết bị (Device Inventory Service):} Mọi máy trạm của nhân viên đều được cấp chứng thư số định danh lưu trong chip bảo mật phần cứng (TPM) và liên tục cập nhật trạng thái bản vá, cấu hình bảo mật (\textit{Device Posture}).
        \item \textit{Cổng ủy nhiệm L7 (Access Proxy):} Toàn bộ ứng dụng nội bộ được đưa ra sau lớp Access Proxy kết hợp xác thực SSO và khóa bảo mật phần cứng chống lừa đảo (\textit{Phishing-resistant MFA}).
        \item \textit{Động cơ kiểm soát truy cập (Access Control Engine):} Đánh giá đồng thời định danh người dùng, độ tin cậy thiết bị và phân loại độ nhạy cảm của ứng dụng đích trước khi cấp quyền cho từng phiên HTTP/HTTPS hoặc SSH.
    \end{enumerate}
    \item \textbf{Kết quả đạt được:} Hơn 100.000 nhân viên của Google có thể làm việc từ bất kỳ mạng Internet công cộng nào (quán cà phê, sân bay, tại nhà) với trải nghiệm truy cập ứng dụng tương đương mạng văn phòng mà không cần bật VPN, đồng thời cô lập hoàn toàn hạ tầng ứng dụng khỏi các cuộc tấn công quét mạng Tầng 3.
\end{itemize}

\subsection{Quản lý kết nối trong môi trường Đa đám mây: Case Study Tập đoàn Siemens AG}
\label{subsec:5_1_2}

Đối với các tập đoàn đa quốc gia vận hành song song Trung tâm dữ liệu nội bộ (\textit{On-premises}) và nhiều nền tảng đám mây công cộng (AWS, Microsoft Azure), việc duy trì mạng đường trục MPLS và các đường hầm Site-to-Site VPN tạo ra cấu trúc định tuyến phức tạp. Quá trình chuyển đổi hạ tầng của tập đoàn công nghiệp \textbf{Siemens AG} là minh chứng thực tiễn cho việc ứng dụng ZTNA ở quy mô toàn cầu \cite{gartner2022}.

\begin{itemize}
    \item \textbf{Bối cảnh và vấn đề gặp phải:} Siemens vận hành hạ tầng phục vụ hơn 300.000 nhân viên tại 190 quốc gia. Trước khi chuyển đổi, lưu lượng truy cập từ các chi nhánh và người dùng từ xa đều phải kéo về các trung tâm dữ liệu tập trung qua đường truyền MPLS và VPN (hiện tượng \textit{Traffic Hairpinning}), gây độ trễ cao khi truy cập ứng dụng đám mây và tốn kém chi phí thuê kênh riêng. Đồng thời, hàng trăm địa chỉ IP công cộng của các cổng VPN và tường lửa biên bị lộ diện trên Internet, làm gia tăng bề mặt tấn công.
    \item \textbf{Mô hình kiến trúc ZTNA áp dụng:} Siemens triển khai nền tảng ZTNA gốc đám mây (Zscaler Private Access) theo mô hình phân bổ lai (\textit{Hybrid Connectors}):
    \begin{enumerate}
        \item Đặt các bộ kết nối nhẹ (\textit{App Connectors}) bên trong từng phân vùng mạng ảo (VPC/VNet) trên AWS, Azure và trung tâm dữ liệu nội bộ.
        \item Đóng toàn bộ các cổng kết nối đi vào (\textit{Inbound Ports}) trên tường lửa doanh nghiệp; các App Connector chỉ thiết lập kết nối đi ra (\textit{Outbound-only}) qua cổng 443 tới đám mây ZTNA.
        \item Định tuyến lưu lượng người dùng trực tiếp qua các nút Edge đám mây gần nhất tới đúng ứng dụng đích dựa trên tên miền FQDN.
    \end{enumerate}
    \item \textbf{Kết quả đạt được:} Siemens loại bỏ hoàn toàn các cổng dịch vụ lộ diện công cộng của ứng dụng nội bộ (\textit{Dark Cloud}), triệt tiêu nguy cơ lây lan mã độc giữa các phân vùng đám mây và giảm đáng kể chi phí vận hành đường truyền MPLS truyền thống.
\end{itemize}

\subsection{Cấp quyền cho Nhà thầu và Chống chiếm đoạt tài khoản: Case Study Cloudflare}
\label{subsec:5_1_3}

Việc cấp quyền truy cập cho đối tác bên ngoài, nhà thầu hoặc nhân viên sử dụng thiết bị cá nhân (\textit{BYOD}) luôn tiềm ẩn nguy cơ rò rỉ thông tin xác thực. Quá trình triển khai nội bộ giải pháp \textbf{Cloudflare Access} tại chính công ty Cloudflare đã chứng minh hiệu quả thực tế của mô hình ZTNA trước các cuộc tấn công lừa đảo tinh vi \cite{cloudflare2022}.

\begin{itemize}
    \item \textbf{Bối cảnh và vấn đề gặp phải:} Đội ngũ kỹ sư và các nhà thầu đối tác của Cloudflare phân tán toàn cầu, thường xuyên cần truy cập vào các hệ thống phát triển nội bộ nhạy cảm (kho mã nguồn Git, Jira, Jenkins, máy chủ SSH). Hệ thống VPN cũ gặp khó khăn trong việc phân tách quyền chi tiết cho từng nhà thầu thời vụ và không thể ngăn chặn kẻ tấn công xâm nhập nếu thông tin đăng nhập bị đánh cắp qua Phishing.
    \item \textbf{Mô hình kiến trúc ZTNA áp dụng:} Cloudflare thay thế VPN bằng lớp Ủy nhiệm ngược nhận biết định danh (\textit{Identity-Aware Reverse Proxy}) kết hợp hai cơ chế bảo vệ cốt lõi:
    \begin{enumerate}
        \item \textit{Truy cập không tác nhân qua trình duyệt (Agentless Web/HTML5 Isolation):} Nhà thầu và đối tác truy cập ứng dụng Web hoặc giao diện dòng lệnh SSH trực tiếp qua trình duyệt được kiểm soát chính sách, không cần cài đặt phần mềm VPN trên máy cá nhân và chỉ được cấp quyền tạm thời theo từng dự án.
        \item \textit{Xác thực phần cứng FIDO2 kết hợp mTLS:} Mọi phiên truy cập qua ZTNA Proxy đều bắt buộc xác thực bằng khóa bảo mật vật lý chuẩn FIDO2/WebAuthn (YubiKey) gắn liền với chứng thư số trên thiết bị được công ty cấp phát.
    \end{enumerate}
    \item \textbf{Kết quả thực tế (Ngăn chặn chiến dịch tấn công Oktapus năm 2022):} Vào tháng 7/2022, nhóm tin tặc thực hiện chiến dịch Phishing quy mô lớn (\textit{Oktapus}) nhắm vào hơn 130 tổ chức công nghệ thông qua tin nhắn SMS giả mạo trang đăng nhập SSO. Mặc dù 3 nhân viên của Cloudflare đã nhập tên đăng nhập và mật khẩu vào trang giả mạo, kẻ tấn công vẫn bị chặn đứng hoàn toàn ở lớp ZTNA Access Proxy do không sở hữu khóa mật mã phần cứng FIDO2 và chứng thư định danh thiết bị hợp lệ, bảo vệ an toàn tuyệt đối cho hệ thống nội bộ \cite{cloudflare2022}.
\end{itemize}

\begin{table}[htbp]
\centering
\caption{Tổng hợp đối chiếu các nghiên cứu tình huống triển khai ZTNA thực tế}
\label{tab:ztna_use_cases}
\small
\begin{tabular}{|>{\raggedright\arraybackslash}p{2.9cm}|>{\raggedright\arraybackslash}p{3.2cm}|>{\raggedright\arraybackslash}p{3.4cm}|>{\raggedright\arraybackslash}p{3.8cm}|}
\hline
\textbf{Case Study thực tế} & \textbf{Kịch bản \& Quy mô} & \textbf{Mô hình ZTNA áp dụng} & \textbf{Hiệu quả bảo mật trọng tâm} \\ \hline
\textbf{Google BeyondCorp} & Lao động kết hợp (>100.000 nhân viên) & Client-initiated (Access Proxy + Device Inventory) & Xóa bỏ hoàn toàn VPN, cô lập mạng L3, xác thực theo ngữ cảnh thiết bị. \\ \hline
\textbf{Siemens AG} & Hạ tầng Đa đám mây (>300.000 nhân viên) & Hybrid Cloud-Broker (Outbound Connectors) & Đóng toàn bộ cổng Inbound, giảm trễ định tuyến, thay thế kênh MPLS/VPN. \\ \hline
\textbf{Cloudflare} & Nhà thầu, BYOD \& Chống Phishing nội bộ & Service-initiated Reverse Proxy + FIDO2 mTLS & Cấp quyền vi mô cho đối tác, chặn đứng thực tế chiến dịch Phishing Oktapus. \\ \hline
\end{tabular}
\end{table}

% ==========================================
% SECTION 5.2
% ==========================================
\section{Thách thức và Rào cản kỹ thuật khi triển khai}
\label{sec:5_2}

\subsection{Vấn đề tương thích với các ứng dụng cũ và giao thức phi web (Legacy Applications)}
\label{subsec:5_2_1}

Một trong những rào cản kỹ thuật lớn nhất khi chuyển đổi từ VPN sang ZTNA là sự hiện diện của các hệ thống ứng dụng kế thừa (\textit{Legacy Applications}) được thiết kế từ trước kỷ nguyên điện toán đám mây:

\begin{itemize}
    \item \textbf{Hạn chế của mô hình Reverse Proxy với giao thức phi HTTP:} Các ứng dụng doanh nghiệp truyền thống (như hệ thống ERP/SAP cũ, phần mềm kế toán nội bộ, hệ thống tổng đài VoIP/SIP hoặc giao thức truyền tệp FTP chủ động) thường vận hành trên các cổng TCP/UDP tùy biến, sử dụng dải cổng động (\textit{Dynamic RPC Ports}) hoặc nhúng cứng địa chỉ IP trong mã nguồn thay vì sử dụng tên miền FQDN. Mô hình ZTNA không tác nhân (\textit{Agentless}) vốn chỉ hỗ trợ HTTP/HTTPS sẽ không thể xử lý các lưu lượng này.
    \item \textbf{Luồng kết nối khởi tạo từ máy chủ (\textit{Server-initiated Flows}):} Kiến trúc ZTNA tiêu chuẩn chỉ hỗ trợ luồng giao tiếp một chiều khởi tạo từ phía người dùng (\textit{Client-to-Server}). Các hệ thống giám sát hạ tầng, đẩy bản vá tự động hoặc dịch vụ gọi ngược từ máy chủ về máy trạm (\textit{Server-to-Client Callbacks}) đòi hỏi doanh nghiệp phải triển khai các tác nhân ZTNA hỗ trợ định tuyến hai chiều ở Tầng 4 hoặc duy trì phân đoạn mạng chuyển tiếp tạm thời trong giai đoạn quá độ.
\end{itemize}

\subsection{Độ phức tạp khi tích hợp với hạ tầng định danh hiện hữu (Active Directory, LDAP, IAM)}
\label{subsec:5_2_2}

ZTNA lấy định danh làm ranh giới bảo mật cốt lõi; do đó, hiệu quả của ZTNA phụ thuộc trực tiếp vào mức độ trưởng thành của hệ thống Quản lý định danh và Truy cập (\textit{IAM}) sẵn có:

\begin{itemize}
    \item \textbf{Phân mảnh kho định danh (\textit{Identity Silos}):} Nhiều tổ chức duy trì đồng thời hệ thống Active Directory (AD) cục bộ cho các ứng dụng nội bộ, thư mục LDAP cho máy chủ Linux và các nhà cung cấp định danh đám mây (như Microsoft Entra ID, Okta) cho ứng dụng SaaS. Việc đồng bộ hóa thuộc tính người dùng, nhóm quyền và trạng thái vô hiệu hóa tài khoản theo thời gian thực giữa các hệ thống này đòi hỏi công sức tích hợp phức tạp qua các giao thức SAML 2.0, OIDC và SCIM.
    \item \textbf{Thiếu hụt dữ liệu đo lường từ thiết bị (\textit{Telemetry Gaps}):} Để động cơ đánh giá tin cậy (\textit{Trust Engine}) đưa ra quyết định chính xác, ZTNA Controller phải kết nối API với các hệ thống quản lý thiết bị (MDM/UEM) và giải pháp phát hiện phản hồi điểm cuối (EDR). Nếu doanh nghiệp chưa chuẩn hóa việc định danh thiết bị bằng chứng thư số PKI, hệ thống ZTNA sẽ gặp khó khăn trong việc phân biệt thiết bị hợp lệ với thiết bị giả mạo.
\end{itemize}

\subsection{Rào cản văn hóa tổ chức, quy trình vận hành và thói quen người dùng}
\label{subsec:5_2_3}

Bên cạnh các yếu tố công nghệ, yếu tố con người và cơ cấu tổ chức thường quyết định sự thành bại của dự án chuyển đổi Zero Trust:

\begin{itemize}
    \item \textbf{Sự phân tách giữa các bộ phận kỹ thuật (\textit{Organizational Silos}):} Trong mô hình truyền thống, đội ngũ Mạng (\textit{Network Operations}) quản lý tường lửa và VPN, đội ngũ An toàn thông tin (\textit{SecOps}) giám sát cảnh báo, còn đội ngũ Phát triển (\textit{App Owners}) quản lý ứng dụng. Việc triển khai ZTNA buộc các bộ phận này phải phối hợp chặt chẽ để lập bản đồ luồng dữ liệu ứng dụng và xây dựng chính sách truy cập chi tiết đến từng chức năng.
    \item \textbf{Ma sát trải nghiệm và tâm lý kháng cự của người dùng:} Nếu các ngưỡng đánh giá tin cậy động được cấu hình quá khắt khe (ví dụ: yêu cầu xác thực MFA lặp lại liên tục khi người dùng chuyển đổi mạng Wi-Fi hoặc cập nhật trình duyệt), nhân viên sẽ nảy sinh tâm lý mệt mỏi xác thực (\textit{MFA Fatigue}) và tìm cách lách quy trình bảo mật. Doanh nghiệp cần kết hợp truyền thông nội bộ và thiết lập chính sách chuyển đổi theo từng giai đoạn để giảm thiểu gián đoạn công việc.
\end{itemize}

% ==========================================
% SECTION 5.3
% ==========================================
\section{Lộ trình 5 bước chuyển đổi sang Kiến trúc Zero Trust (ZTA)}
\label{sec:5_3}

Để tránh rủi ro gián đoạn hoạt động kinh doanh khi thay thế toàn bộ hạ tầng mạng cùng lúc, doanh nghiệp cần áp dụng phương pháp luận triển khai cuốn chiếu gồm 5 bước chuẩn hóa:

\subsection{Bước 1: Định danh dữ liệu và tài sản quan trọng (Define the Protect Surface)}
\label{subsec:5_3_1}

Thay vì cố gắng bảo vệ toàn bộ bề mặt tấn công (\textit{Attack Surface}) rộng lớn và luôn biến đổi của toàn hệ thống, doanh nghiệp cần thu hẹp phạm vi tập trung vào các \textbf{Bề mặt cần bảo vệ (\textit{Protect Surface})}. Mỗi Protect Surface bao gồm các thành phần cốt lõi thuộc nhóm \textbf{DAAS}:
\begin{itemize}
    \item \textbf{Dữ liệu (\textit{Data}):} Dữ liệu định danh khách hàng (PII), thông tin thẻ thanh toán (PCI), sở hữu trí tuệ, bí mật kinh doanh.
    \item \textbf{Ứng dụng (\textit{Applications}):} Các phần mềm lõi phục vụ vận hành như ERP, CRM, Core Banking, hệ thống quản lý nhân sự.
    \item \textbf{Tài sản (\textit{Assets}):} Hệ thống máy chủ cơ sở dữ liệu, thiết bị điều khiển công nghiệp (SCADA/IoT), điểm bán hàng.
    \item \textbf{Dịch vụ (\textit{Services}):} Các dịch vụ hạ tầng nền tảng như DNS, DHCP, Active Directory, dịch vụ NTP.
\end{itemize}
Mỗi Protect Surface có quy mô nhỏ, xác định được và hoàn toàn nằm trong tầm kiểm soát của tổ chức, giúp việc thiết lập ranh giới bảo mật trở nên khả thi.

\subsection{Bước 2: Vẽ bản đồ luồng giao tiếp dữ liệu (Map the Transaction Flows)}
\label{subsec:5_3_2}

Sau khi xác định được Protect Surface ưu tiên, đội ngũ kỹ thuật tiến hành phân tích luồng lưu lượng mạng thực tế để hiểu rõ cách thức các thành phần tương tác với nhau:
\begin{itemize}
    \item Thu thập nhật ký lưu lượng từ tường lửa, bộ chuyển mạch, hệ thống giám sát mạng (NetFlow) và nhật ký xác thực ứng dụng.
    \item Xác định chính xác nhóm người dùng nào, thiết bị nào, dịch vụ phụ trợ nào đang giao tiếp với tài sản trong Protect Surface, sử dụng giao thức gì và tần suất ra sao. Kết quả của bước này là sơ đồ phụ thuộc ứng dụng (\textit{Application Dependency Mapping}), giúp loại bỏ các luồng kết nối dư thừa không phục vụ nghiệp vụ.
\end{itemize}

\subsection{Bước 3: Thiết kế kiến trúc mạng Zero Trust xung quanh tài sản (Architect a Zero Trust Network)}
\label{subsec:5_3_3}

Dựa trên bản đồ luồng giao tiếp đã lập, kiến trúc sư hệ thống bố trí các điểm thực thi chính sách (\textit{Policy Enforcement Points - PEP}) bao bọc sát xung quanh từng Protect Surface để tạo thành các vi chu vi (\textit{Micro-perimeters}):
\begin{itemize}
    \item Triển khai ZTNA Gateway hoặc ZTNA Connector ngay phía trước máy chủ ứng dụng, đóng toàn bộ các cổng kết nối trực tiếp từ bên ngoài.
    \item Tích hợp mặt phẳng kiểm soát (ZTNA Controller) với hệ thống quản lý định danh tập trung (IdP/SSO) và các nguồn dữ liệu ngữ cảnh thiết bị (EDR/MDM) để sẵn sàng thực thi cơ chế xác thực trước khi kết nối.
\end{itemize}

\subsection{Bước 4: Xây dựng tập chính sách Zero Trust (Create the Zero Trust Policy)}
\label{subsec:5_3_4}

Ở bước này, các quy tắc kiểm soát truy cập chi tiết dựa trên thuộc tính (ABAC) được thiết lập theo nguyên tắc đặc quyền tối thiểu (\textit{Least Privilege}). Để đảm bảo chính sách không bỏ sót ngữ cảnh, tập quy tắc được xây dựng dựa trên \textbf{Phương pháp luận Kipling} gồm 6 tham số:
\begin{enumerate}
    \item \textbf{Ai (\textit{Who}):} Định danh người dùng hoặc nhóm vai trò cụ thể được phép truy cập (xác thực qua IdP).
    \item \textbf{Cái gì (\textit{What}):} Ứng dụng, API hoặc phân vùng dữ liệu đích cụ thể trong Protect Surface.
    \item \textbf{Khi nào (\textit{When}):} Khung thời gian hoặc thời hạn phiên làm việc được phép truy cập.
    \item \textbf{Ở đâu (\textit{Where}):} Vị trí địa lý, phân vùng mạng nguồn hoặc trạng thái bảo mật của thiết bị đầu cuối (\textit{Device Posture}).
    \item \textbf{Tại sao (\textit{Why}):} Mục đích nghiệp vụ hợp lệ giải trình cho nhu cầu truy cập tài nguyên nhạy cảm.
    \item \textbf{Như thế nào (\textit{How}):} Giao thức truyền tải bắt buộc (ví dụ: chỉ cho phép HTTPS với mTLS, yêu cầu kiểm tra mã độc thời gian thực).
\end{enumerate}
Mọi lưu lượng không khớp với tập chính sách Kipling đã định nghĩa đều bị chặn mặc định (\textit{Implicit Deny}).

\subsection{Bước 5: Giám sát, phân tích và cải tiến liên tục (Monitor and Maintain)}
\label{subsec:5_3_5}

Triển khai Zero Trust là một chu trình lặp liên tục chứ không phải một dự án tĩnh kết thúc sau khi cài đặt:
\begin{itemize}
    \item Toàn bộ nhật ký truy cập Tầng 7, sự kiện cấp/thu hồi phiên và dữ liệu viễn trắc thiết bị từ hệ thống ZTNA được đẩy tập trung về nền tảng SIEM/SOAR.
    \item Đội ngũ vận hành an ninh thường xuyên rà soát các chính sách truy cập không còn sử dụng (\textit{Unused Policies}), tinh chỉnh ngưỡng điểm tin cậy, sau đó lặp lại quy trình từ Bước 1 cho Protect Surface tiếp theo cho đến khi toàn bộ hạ tầng doanh nghiệp được chuyển đổi hoàn toàn.
\end{itemize}

% ==========================================
% SECTION 5.4
% ==========================================
\section{Xu hướng hội tụ công nghệ tương lai}
\label{sec:5_4}

\subsection{Khung kiến trúc SASE (Secure Access Service Edge) và SSE (Security Service Edge)}
\label{subsec:5_4_1}

Sự phân tán của người dùng và ứng dụng đám mây thúc đẩy xu hướng hợp nhất hạ tầng mạng diện rộng và các dịch vụ bảo mật thành một nền tảng đám mây duy nhất, được định nghĩa thông qua hai khung kiến trúc chiến lược:

\begin{itemize}
    \item \textbf{Secure Access Service Edge (SASE):} Là kiến trúc hội tụ giữa công nghệ mạng định nghĩa bằng phần mềm (\textbf{SD-WAN}) và bộ dịch vụ bảo mật gốc đám mây (\textbf{SSE}), được phân phối trực tiếp từ mạng lưới các nút Edge phân tán toàn cầu. SASE cho phép doanh nghiệp định tuyến lưu lượng tối ưu đồng thời áp dụng chính sách bảo mật nhất quán tại điểm gần người dùng nhất.
    \item \textbf{Security Service Edge (SSE):} Đóng vai trò là trụ cột an toàn thông tin của kiến trúc SASE, tập trung vào việc hợp nhất toàn bộ các năng lực kiểm soát truy cập và bảo vệ dữ liệu đám mây trên một mặt phẳng quản trị thống nhất, bao gồm ba thành phần nòng cốt: \textbf{ZTNA}, \textbf{CASB} và \textbf{SWG}.
\end{itemize}

\subsection{Tích hợp hợp nhất ZTNA với CASB và SWG}
\label{subsec:5_4_2}

Trong hệ sinh thái SSE hiện đại, ZTNA không hoạt động đơn lẻ mà tích hợp chặt chẽ với hai công nghệ bảo mật biên đám mây khác để bao phủ toàn diện mọi hướng lưu lượng của người dùng:

\begin{itemize}
    \item \textbf{ZTNA (Kiểm soát ứng dụng riêng tư):} Đảm nhiệm việc bảo vệ và cấp quyền truy cập vi mô tới các ứng dụng nội bộ đặt tại trung tâm dữ liệu riêng hoặc hạ tầng đám mây cá nhân (Private Cloud/IaaS).
    \item \textbf{CASB (\textit{Cloud Access Security Broker}):} Kiểm soát các ứng dụng đám mây công cộng được doanh nghiệp cấp phép (Sanctioned SaaS như Microsoft 365, Salesforce, Google Workspace), phát hiện ứng dụng đám mây trái phép (\textit{Shadow IT}) và thực thi chống thất thoát dữ liệu (DLP) ở cấp độ API.
    \item \textbf{SWG (\textit{Secure Web Gateway}):} Đảm nhiệm lọc lưu lượng truy cập Internet công cộng, phân loại URL độc hại, giải mã kiểm tra lưu lượng TLS và ngăn chặn mã độc từ các trang web bên ngoài.
\end{itemize}

Sự hợp nhất này giúp người dùng chỉ cần cài đặt một tác nhân hợp nhất duy nhất (\textit{Unified Agent}) trên thiết bị đầu cuối, trong khi quản trị viên vận hành một bộ chính sách bảo mật đồng nhất cho cả ứng dụng nội bộ, ứng dụng SaaS và truy cập Web.

\begin{figure}[htbp]
\centering
\resizebox{\textwidth}{!}{%
\begin{tikzpicture}[font=\small]
    % Cột trái: Người dùng & Chi nhánh
    \node[draw, rounded corners=4pt, fill=blue!8, minimum width=3.3cm, minimum height=1.15cm, align=center] (user1) at (1.65, 3.8) {\textbf{Người dùng từ xa}\\(Remote / Hybrid)};
    \node[draw, rounded corners=4pt, fill=blue!8, minimum width=3.3cm, minimum height=1.15cm, align=center] (user2) at (1.65, 2.0) {\textbf{Văn phòng Chi nhánh}\\(Branch Offices)};
    \node[draw, rounded corners=4pt, fill=blue!8, minimum width=3.3cm, minimum height=1.15cm, align=center] (user3) at (1.65, 0.2) {\textbf{Đối tác / BYOD}\\(Third-party)};

    % Khung trung tâm: SASE Cloud Edge (Mở rộng chiều ngang lên 9.4cm)
    \draw[draw=teal!80!black, thick, rounded corners=6pt, fill=teal!4] (4.1, -0.65) rectangle (13.5, 4.95);
    \node[font=\bfseries, text=teal!80!black] at (8.8, 4.6) {NỀN TẢNG HỘI TỤ SASE (CLOUD EDGE PoP)};

    % Lớp SD-WAN
    \draw[draw=cyan!70!black, rounded corners=4pt, fill=cyan!15] (4.4, 3.55) rectangle (13.2, 4.25);
    \node[align=center] at (8.8, 3.9) {\textbf{SD-WAN} (Định tuyến tối ưu \& QoS đa đường truyền)};

    % Khung nét đứt SSE
    \draw[draw=violet!70!black, dashed, rounded corners=4pt, fill=white] (4.4, -0.4) rectangle (13.2, 3.35);
    \node[font=\bfseries\scriptsize, text=black!80] at (8.8, 3.08) {SECURITY SERVICE EDGE (SSE)};

    % 4 thành phần bên trong SSE
    \node[draw, rounded corners=3pt, fill=green!15, minimum width=3.9cm, minimum height=1.1cm, align=center] (ztna) at (6.65, 2.1) {\textbf{ZTNA}\\Ứng dụng nội bộ};
    \node[draw, rounded corners=3pt, fill=orange!15, minimum width=3.9cm, minimum height=1.1cm, align=center] (casb) at (10.95, 2.1) {\textbf{CASB}\\Ứng dụng SaaS};
    \node[draw, rounded corners=3pt, fill=purple!12, minimum width=3.9cm, minimum height=1.1cm, align=center] (swg) at (6.65, 0.45) {\textbf{SWG}\\Lọc Web / Internet};
    \node[draw, rounded corners=3pt, fill=red!12, minimum width=3.9cm, minimum height=1.1cm, align=center] (ai) at (10.95, 0.45) {\textbf{AI / UEBA}\\Chính sách động};

    % Cột phải: Đích đến
    \node[draw, rounded corners=4pt, fill=gray!15, minimum width=3.3cm, minimum height=1.15cm, align=center] (dest1) at (15.95, 3.8) {\textbf{Data Center /}\\Private Cloud};
    \node[draw, rounded corners=4pt, fill=gray!15, minimum width=3.3cm, minimum height=1.15cm, align=center] (dest2) at (15.95, 2.0) {\textbf{Ứng dụng SaaS}\\(M365, Salesforce)};
    \node[draw, rounded corners=4pt, fill=gray!15, minimum width=3.3cm, minimum height=1.15cm, align=center] (dest3) at (15.95, 0.2) {\textbf{Mạng Internet}\\Công cộng};

    % Mũi tên trái -> SASE
    \draw[->, thick] (user1.east) -- (4.1, 3.8);
    \draw[->, thick] (user2.east) -- (4.1, 2.0);
    \draw[->, thick] (user3.east) -- (4.1, 0.2);

    % Mũi tên SASE -> phải
    \draw[->, thick] (13.5, 3.8) -- (dest1.west);
    \draw[->, thick] (13.5, 2.0) -- (dest2.west);
    \draw[->, thick] (13.5, 0.2) -- (dest3.west);
\end{tikzpicture}%
}
\caption{Mô hình kiến trúc hội tụ SASE và hợp nhất các thành phần bảo mật SSE}
\label{fig:sase_sse_convergence}
\end{figure}

\subsection{Ứng dụng AI/ML và Phân tích hành vi người dùng (UEBA) trong điều chỉnh chính sách động}
\label{subsec:5_4_3}

Trong các thế hệ giải pháp ZTNA tiếp theo, trí tuệ nhân tạo (AI), học máy (ML) và phân tích hành vi người dùng và thực thể (\textbf{UEBA} - \textit{User and Entity Behavior Analytics}) được tích hợp sâu vào động cơ đánh giá tin cậy nhằm chuyển dịch từ kiểm soát dựa trên quy tắc tĩnh sang phòng thủ thích ứng tự động:

\begin{itemize}
    \item \textbf{Xây dựng hồ sơ hành vi cơ sở (\textit{Behavioral Baselining}):} Các mô hình học máy không giám sát phân tích dữ liệu lịch sử về khung giờ làm việc, vị trí địa lý đăng nhập, thiết bị quen dùng, tốc độ gõ phím và tần suất truy xuất dữ liệu của từng người dùng để thiết lập hồ sơ hành vi chuẩn mực.
    \item \textbf{Phát hiện bất thường theo thời gian thực:} Hệ thống UEBA liên tục đối chiếu hành vi hiện tại với đường cơ sở để nhận diện các chỉ báo xâm nhập tinh vi mà quy tắc truyền thống bỏ sót, chẳng hạn như hiện tượng di chuyển bất khả thi (\textit{Impossible Travel}), truy xuất hàng loạt tệp tin nhạy cảm ngoài giờ hành chính hoặc chuỗi gọi API bất thường từ một tài khoản đã xác thực hợp lệ.
    \item \textbf{Tự động hóa phản ứng chính sách (\textit{Autonomous Policy Adaptation}):} Khi thuật toán phát hiện mức độ rủi ro tăng đột biến, điểm tin cậy $T_{\text{score}}$ lập tức bị hạ thấp. Hệ thống tự động kích hoạt quy trình phản ứng theo cấp độ: yêu cầu xác thực sinh trắc học bổ sung, giới hạn quyền chỉ đọc (\textit{Read-only}), cô lập phiên vào môi trường trình duyệt ảo từ xa (\textit{Remote Browser Isolation - RBI}) hoặc chấm dứt hoàn toàn phiên kết nối mà không cần chờ quản trị viên can thiệp thủ công.
\end{itemize}